\section{Reviewer 1}

\subsection{Major Comments}

\begin{enumerate}
	\item The research objective of this paper should be clarified. I can see four different aims that are addressed at least partially in this paper. The first is to develop the best performing classifier to predict battles in a civil conflict. The second is to understand whether a latent variable network model performs better than GLM in predicting battles, with the intent of interpreting the performance gain as the benefit of accounting for network dependencies. The third is to develop an explanatory model of battles from which we can learn about the factors that serve as indicators of civil conflict. The first is to propose and illustrate the application of a latent variable network model that can incorporate exogenous variation in the vertex set. Once the authors clarify their objective, it would be useful to provide additional analysis that builds toward their aim. For example, if the goal is to develop an effective classifier, it would be valuable to see comparisons to random forests, neural networks, and other learning approaches at are designed to optimize predictive performance.
	\begin{itemize}
		\item \textcolor{blue}{ \emph{
		Insert great response.
		}}
	\end{itemize}	
	\item Second, the empirical modeling presented in this paper is generally very carefully executed. One aspect of the model specification that I would recommend exploring is the degree to which there is temporal heterogeneity in the parameter values. I am not an expert in civil conflict, but it seems likely that the factors governing battle instances will shift over time throughout the conflict.
	\begin{itemize}
		\item \textcolor{blue}{ \emph{
		Insert great response.
		}}
	\end{itemize}
	\item I recommend that the authors add the TERGM as a comparison model, or provide additional justification for not comparing their model to TERGM. For example, perhaps after the goals of the paper are clarified, it would be clear that there would be no need to compare their model to yet another network model that can account for dependence ... it would be great to see a robustness check in which the authors evaluate whether the effects of covariates are robust to modeling with TERGM.
	\begin{itemize}
		\item \textcolor{blue}{ \emph{
		Insert great response.
		}}
	\end{itemize}	
\end{enumerate}

\subsection{Minor Comments}

\begin{enumerate}
	\item I would like to see a discussion of setting the priors used for inference with this model. In particular, how does the specification of the priors relate to the predictive performance of the latent variable model? If the priors make any sort of difference, how should the priors be set to assure that the model performs as well as possible while also assuring that the model is not trained on the test data.
	\begin{itemize}
		\item \textcolor{blue}{ \emph{
		We discuss the priors on page ...
		}}
	\end{itemize}	
	\item Second, for the predictive experiment, it would be great to see if the comparison of the AME and GLM is sensitive to k---the number of groups into which the data is split ... it would be great to see if the predictive performance of AME relative to GLM varied between, say k=5 and k=20.
	\begin{itemize}
		\item \textcolor{blue}{ \emph{
		We insert a mention of the value of $k$ used in our analysis on page X. We also, in the appendix, compare the predictive value for values of $k$ ranging from 5 to 20.
		}}
	\end{itemize}	
\end{enumerate}